\documentclass[../main.tex]{subfiles}

\begin{document}

\section{Kiến trúc tổng thể hệ thống}
Hệ thống TryHairStyle được tổ chức theo kiến trúc client--server với bốn thành phần chính, đảm bảo khả năng mở rộng và xử lý hiệu quả các tác vụ tính toán nặng:

\begin{enumerate}
    \item \textbf{Frontend (React/Vite):}
    Giao diện web cho phép người dùng tải ảnh đầu vào, lựa chọn kiểu tóc tham chiếu, tùy chỉnh màu tóc và theo dõi kết quả. Frontend được thiết kế để tương tác thời gian thực với backend thông qua các API.

    \item \textbf{Backend API (FastAPI):}
    Thành phần này tiếp nhận các yêu cầu HTTP từ phía người dùng, thực hiện kiểm tra dữ liệu đầu vào và khởi tạo tác vụ xử lý. Ngoài ra, backend còn chịu trách nhiệm quản lý trạng thái và cung cấp các tài nguyên tĩnh như ảnh kết quả.

    \item \textbf{Worker AI (Celery):}
    Worker đảm nhận việc xử lý các tác vụ suy luận AI. Sau khi nhận nhiệm vụ từ hàng đợi, worker thực thi toàn bộ pipeline bao gồm tiền xử lý, suy luận mô hình và hậu xử lý trên GPU, đảm bảo hiệu suất cho các tác vụ nặng.

    \item \textbf{Message Broker (Redis):}
    Redis đóng vai trò trung gian truyền tải tác vụ giữa backend và worker, đồng thời lưu trữ trạng thái xử lý. Cơ chế hàng đợi giúp hệ thống xử lý bất đồng bộ và tăng khả năng chịu tải.
\end{enumerate}

Kiến trúc này tách biệt rõ ràng giữa luồng tiếp nhận yêu cầu và luồng xử lý suy luận, cho phép backend phản hồi nhanh chóng mà không bị chặn bởi các tác vụ tính toán phức tạp. Nhờ đó, hệ thống hạn chế tình trạng timeout và nâng cao trải nghiệm người dùng.

\subsection{Cơ chế điều kiện kép}
Để vừa bảo toàn khuôn mặt vừa chuyển giao kiểu tóc, hệ thống sử dụng hai nhánh điều kiện hóa song song, mỗi nhánh đảm nhận một vai trò riêng biệt nhưng bổ trợ lẫn nhau:

\begin{itemize}
    \item \textbf{Nhánh bảo toàn danh tính:}
    Ảnh người dùng được phân tích bằng InsightFace để trích xuất embedding khuôn mặt 512 chiều, đóng vai trò biểu diễn danh tính. Embedding này được kết hợp với ControlNet Depth nhằm duy trì cấu trúc hình học, đồng thời tích hợp InstantID để tăng cường khả năng giữ nguyên đặc trưng nhận dạng. Đối với các trường hợp khuôn mặt nghiêng ($|yaw| \geq 45^\circ$), hệ thống tự động chuyển sang sử dụng AdaFace IR\nobreakdash-100 để đảm bảo độ ổn định của embedding.

    \item \textbf{Nhánh truyền tải phong cách tóc:}
    Ảnh tóc tham chiếu được mã hóa thành vector đặc trưng thông qua IP-Adapter Plus trong pipeline chuẩn, hoặc HairTextureEncoder (ResNet\nobreakdash-50) trong pipeline tùy chỉnh. Vector này chứa thông tin về cấu trúc, hình dạng và kết cấu tóc, từ đó dẫn hướng mô hình sinh ảnh tạo ra kiểu tóc phù hợp với tham chiếu.
\end{itemize}

Hai nhánh điều kiện này được tích hợp đồng thời vào quá trình suy luận của mô hình khuếch tán, giúp cân bằng giữa việc giữ nguyên danh tính khuôn mặt và áp dụng chính xác kiểu tóc mong muốn.

\section{Thiết kế các mô-đun xử lý}
Phần backend được chia thành các service module độc lập, mỗi module đảm nhận một chức năng riêng biệt.

\subsection{FaceInfoService -- Phát hiện và phân tích khuôn mặt}
Service này sử dụng pipeline ba tầng:
\begin{enumerate}
    \item \textbf{YOLOv8-Face (TrainingFaceDetector):} Phát hiện toàn bộ khuôn mặt trong ảnh, hỗ trợ cả khuôn mặt nghiêng và một phần bị che. Kết quả qua bước NMS (Non-Maximum Suppression) để loại bỏ bounding box trùng lặp.
    \item \textbf{InsightFace (antelopev2):} Phân tích chi tiết từng khuôn mặt, trích xuất 106 landmarks, embedding ArcFace 512 chiều và keypoints.
    \item \textbf{TrainingEmbedder (AdaFace fallback):} Với khuôn mặt mà InsightFace bỏ sót (profile, góc nghiêng lớn), hệ thống dùng AdaFace IR\nobreakdash-100 kết hợp MTCNN alignment để trích xuất embedding. Ngưỡng chuyển đổi là $|yaw| = 45\degree$.
\end{enumerate}
Ngoài ra, service tích hợp \textbf{TrainingReconstructor3D} (3DDFA V2) để reconstruct dense 3D face mesh. Dữ liệu 3D vertices được dùng để ngoại suy vùng da đầu (scalp projection) trong trường hợp người dùng ít tóc hoặc trọc đầu.

\subsection{SegmentationService -- Phân đoạn ngữ nghĩa}
Service sử dụng \textbf{SegFormer} (jonathandinu/face-parsing) để phân đoạn ảnh thành 19 lớp ngữ nghĩa (skin, hair, hat, nose, eyes, v.v.). Hai nhóm mask chính được tạo:
\begin{itemize}
    \item \textbf{Hair mask:} Gộp lớp tóc (class 13) và mũ/nón (class 14) để inpainting thay thế toàn bộ vùng tóc, kể cả khi người dùng đội nón.
    \item \textbf{Face mask:} Tổng hợp 12 lớp thuộc khuôn mặt (skin, nose, eyes, brows, ears, mouth, lips) để bảo vệ khuôn mặt khỏi bị sinh tóc chồng lên.
\end{itemize}
Mask được dilate để mở rộng biên, đảm bảo vùng tóc được nhận diện đầy đủ.

Tuy nhiên, cơ chế bảo vệ toàn bộ vùng da mặt tạo ra một hạn chế đối với các kiểu tóc có phần mái (bangs) phủ xuống trán. Do vùng trán thuộc lớp skin (class~1) nằm trong tập face mask, mô hình sinh ảnh không thể tạo tóc tại vùng này, dẫn đến kết quả bị cắt ngang bất tự nhiên tại đường viền da trán. Để giải quyết vấn đề này, hệ thống áp dụng kỹ thuật \textbf{Forehead Unmasking} nhằm loại bỏ có chọn lọc phần bảo vệ vùng trán trong face mask mà vẫn duy trì bảo vệ cho các vùng nhạy cảm phía dưới. Quy trình gồm ba bước:
\begin{enumerate}
    \item Xác định tọa độ $y_{min}$ của tập hợp các pixel thuộc lớp mắt và lông mày (class 3--7 trong SegFormer). Giá trị $y_{min}$ này đóng vai trò đường ranh giới nằm ngang, phân tách vùng trán phía trên và vùng mặt dưới chứa các đặc trưng nhận dạng.
    \item Áp dụng biên an toàn bằng cách dịch đường ranh giới lên trên một khoảng $\delta = 0{,}02 \times H$ (với $H$ là chiều cao ảnh) nhằm tránh xâm phạm vùng lông mày.
    \item Gán giá trị 0 cho toàn bộ pixel của face mask có tọa độ $y < y_{min} - \delta$, loại bỏ lớp bảo vệ tại vùng trán trong khi giữ nguyên lớp bảo vệ cho vùng mắt, mũi, miệng, má và cằm.
\end{enumerate}
Nhờ kỹ thuật này, mô hình khuếch tán có thể sinh tóc mái phủ tự nhiên xuống trán mà không ảnh hưởng đến các đặc trưng nhận dạng khuôn mặt phía dưới. Trong trường hợp SegFormer không phát hiện được pixel thuộc lớp mắt hoặc lông mày (do ảnh bị che khuất nặng hoặc góc nghiêng lớn), hệ thống bỏ qua bước Forehead Unmasking và giữ nguyên face mask ban đầu nhằm đảm bảo tính ổn định.

\subsection{Kỹ thuật mặt nạ động}
Trong bài toán thử tóc ảo, tóc tham chiếu có thể dài hoặc phồng hơn đáng kể so với tóc gốc. Hệ thống xây dựng cơ chế \texttt{expand\_hair\_mask} với các bước:
\begin{itemize}
    \item \textbf{Phát hiện trường hợp trọc/ít tóc:} Khi diện tích mask tóc nhỏ hơn 100 pixel, hệ thống tạo scalp mask bằng 3DDFA V2 (ngoại suy đỉnh đầu từ forehead vertices, chiếu convex hull) hoặc fallback sang OpenCV (crop tóc mẫu, scale theo face width, paste lên trán).
    \item \textbf{Mở rộng theo tỷ lệ:} Nếu diện tích tóc tham chiếu $> 1.5 \times$ diện tích tóc user, mask được dilate với kernel hình chữ nhật hướng ngang (tóc phồng), dọc (tóc dài) hoặc vuông (tóc đều), từ 1--8 iterations tỷ lệ với chênh lệch diện tích.
    \item \textbf{Bảo vệ khuôn mặt (3 lớp):} Face mask được dilate thêm 10px tạo buffer, trừ khỏi expanded mask, rồi smooth biên bằng Gaussian blur.
\end{itemize}

\subsection{HairDiffusionService -- Suy luận Diffusion}
Service hỗ trợ hai chế độ pipeline:
\begin{enumerate}
    \item \textbf{Standard Pipeline:} Sử dụng \texttt{StableDiffusionXLControlNetInpaintPipeline} với ControlNet Depth + IP-Adapter Plus. Ảnh tham chiếu được dùng trực tiếp qua IP-Adapter.
    \item \textbf{Custom Pipeline:} Sử dụng UNet 13 kênh đầu vào (noisy latents + masked latents + mask + reference hair latents) được tinh chỉnh bằng LoRA. Style embedding (2048 chiều từ HairTextureEncoder) và identity embedding (512 chiều từ InsightFace) được inject qua CrossAttentionInjector. Quá trình khử nhiễu áp dụng Classifier-Free Guidance.
\end{enumerate}
Tất cả ảnh được resize về 512$\times$512 khi suy luận rồi upscale về kích thước gốc bằng Lanczos.

Sau khi giải mã VAE, bước hậu xử lý hòa trộn mềm tại biên mask bằng Gaussian blur ($21 \times 21$, $\sigma = 10$) giúp chuyển tiếp tự nhiên giữa vùng tóc mới và vùng ảnh gốc.

\subsection{HairColorService -- Đổi màu tóc}
Service đổi màu tóc bằng kỹ thuật HSV Color Transfer mà không cần chạy lại mô hình khuếch tán:
\begin{enumerate}
    \item Chuyển ảnh sang không gian HSV.
    \item Thay thành phần Hue (sắc độ) bằng giá trị màu đích.
    \item Blend Saturation (độ bão hòa): giữ 20\% gốc, thay 80\% theo màu mới.
    \item Điều chỉnh Value (độ sáng) nhẹ theo tỷ lệ sáng/tối của màu đích (giới hạn $0.6$--$1.4$).
    \item Blend bằng Gaussian blur trên biên mask để tránh viền cứng.
\end{enumerate}
Hệ thống cung cấp 15 preset màu (đen, nâu, vàng, đỏ, hồng, xanh dương, tím, v.v.) và hỗ trợ hex code tùy chỉnh.

\section{Thiết kế luồng ứng dụng bất đồng bộ}
Do thời gian suy luận kéo dài vài giây, hệ thống áp dụng mô hình xử lý bất đồng bộ:
\begin{enumerate}
    \item Người dùng gửi yêu cầu qua \texttt{POST /generate}.
    \item Backend lưu ảnh upload, tạo Celery task và trả về \texttt{task\_id} ngay lập tức.
    \item Celery worker nhận task từ Redis, thực hiện pipeline AI, cập nhật trạng thái (\texttt{PROCESSING}) qua mỗi bước (Face Analysis, Creating Mask, Estimating Depth, Generating Hair, v.v.).
    \item Frontend poll \texttt{GET /status/\{task\_id\}} mỗi 2 giây cho đến khi nhận \texttt{SUCCESS} hoặc \texttt{FAILURE}.
    \item Kết quả được lưu theo session folder riêng biệt và phục vụ qua static file server.
\end{enumerate}

\section{Thiết kế API RESTful}
Backend cung cấp các endpoint chính được mô tả trong Bảng~\ref{tab:api-endpoints}.

\begin{table}[H]
    \centering
    \caption{Danh sách các API endpoint chính}
    \label{tab:api-endpoints}
    \begin{tabularx}{\textwidth}{|>{\raggedright\arraybackslash}p{1.5cm}|>{\raggedright\arraybackslash}p{3.5cm}|X|}
        \hline
        \textbf{Method} & \textbf{Endpoint} & \textbf{Mô tả} \\
        \hline
        POST & \texttt{/generate} & Nhận 2 ảnh (chân dung + tóc tham chiếu), prompt, màu tóc tùy chọn. Trả về \texttt{task\_id}. \\
        \hline
        POST & \texttt{/detect-faces} & Quét và cắt chân dung các khuôn mặt trong ảnh. Chạy qua Celery. \\
        \hline
        POST & \texttt{/colorize} & Đổi màu tóc trên ảnh gốc (không cần Diffusion Model). \\
        \hline
        GET & \texttt{/status/\{task\_id\}} & Kiểm tra trạng thái tác vụ (PENDING, PROCESSING, SUCCESS, FAILURE). \\
        \hline
        GET & \texttt{/colors} & Trả về danh sách 15 preset màu tóc. \\
        \hline
        GET & \texttt{/random-pair} & Lấy cặp ảnh mẫu ngẫu nhiên từ dataset FFHQ. \\
        \hline
    \end{tabularx}
\end{table}

\end{document}
